Dimensional maps of psychopathology are only as trustworthy as the measurement
that produces them, yet they are typically read as fixed point-coordinates that
hide how sharply each patient is located, how much each instrument contributed,
and how cheaply the map could be measured better. We present the transdiagnostic
dimensional map as a \emph{measurement instrument} rather than a finding. A
missingness-native Bayesian bifactor / exploratory structural-equation model ---
fit to each patient's observed cells only (no imputation), under mixed
likelihoods and a Gaussian-copula continuous block --- returns for every one of
$N=9{,}013$ patients (bipolar disorder, schizophrenia, major depression; 21
sites) a posterior coordinate \emph{and} its uncertainty ellipsoid, and its
closed-form structure makes four instrument properties computable on that one
object. \textbf{Projection:} an expected-a-posteriori score map, reduced to the
eight-dimensional factor space by a Woodbury identity, returns the prior --- never
a false-precise point --- on unmeasured axes. \textbf{Information:} because
indicators are conditionally independent given the factors, posterior precision
accumulates additively, and the quantity that governs localization is the Fisher
information each indicator contributes at its own likelihood family, \emph{not}
its loading --- high-loading, low-prevalence flags contribute almost nothing.
\textbf{Adaptive sampling:} a loading-dependent minimal-indicator rule recovers
most of each axis's reliable signal in three to six items and exposes two axes
that are intrinsically capped by the instrument bank (mania/activation $0.41$,
substance $0.43$) rather than by missingness; a fully patient-adaptive selection
order coincides with the fixed order, showing the population-mean reliability
curve is a \emph{tight} upper bound. \textbf{Value of information:} the entropy
reduction from a not-yet-measured indicator is closed-form (matrix-determinant
lemma), so a greedy rule assembles a shared 27-item battery at mean reliability
$0.70$ and localizes where each cohort is under-measured. We show by simulation
that the posterior coordinate is well-calibrated (empirical coverage $0.95$ matches
nominal, flat across observed-item counts and including the bank-limited axes) and
that its calibration degrades gracefully rather than collapsing under four
violations of the model's assumptions; that the shrinkage-corrected posterior-mean
coordinate improves downstream prognosis over a naive sum-score, most under sparse
measurement (while the posterior's uncertainty does not, on its own, triage who is
predictable for a coarse remission endpoint); and that ranking items by information
rather than loading changes which battery you would build ($+0.11$ reliability at
$27$ items). A five-archetype convex-blend simplex
summarizes the continuum --- which has no natural clusters --- as graded membership
in five named clinical poles, locating a patient $\approx 9.7\times$ more tightly
than the DSM-5 label; we show its modest reconstruction fidelity ($59\%$ of
coordinate variance) is the wrong yardstick for it, since the near-isotropic
variance lets even a random five-dimensional subspace retain more, and the simplex
earns its place by interpretability, not variance retained. The contribution is the
instrument: a map that reports its own uncertainty, weighs instruments by the
information they carry, and prescribes how to measure the next cohort more
cheaply. The clinical and biological findings this map supports are reported in
a companion paper; here the map's biology serves only as a worked example of the
instrument.
